\section{Introduction}\label{introduction}

Retrieval-augmented generation (RAG) systems depend on retrieval quality to ground generated answers, so embedding-model choice directly affects downstream reliability. In operational settings, this choice is rarely made on generic benchmark corpora alone; it is made for local document collections whose language, domain vocabulary, document structure, and latency constraints may differ from public evaluation settings. A model that performs strongly in one context may therefore be less suitable when the deployment corpus changes. \citep{Lewis_RAG_2021}

Public leaderboards and informal local trials remain useful starting points, but they are insufficient as a complete basis for deployment decisions. Generic rankings do not reveal whether a model remains effective on a practitioner's own corpus, while ad hoc testing can confound model behavior with changes in preprocessing, query construction, relevance judgments, retrieval depth, or reporting criteria. Existing local evaluations also tend to treat key concerns separately: they may compare retrieval quality, report operational cost, or preserve reproducible artifacts, but they do not always combine shared-protocol model comparison, reusable benchmark evidence, language and domain comparability, quality--speed reporting, consistency diagnostics, and statistical decision support in one workflow. This gap motivates a controlled local benchmarking process that produces corpus-specific evidence for fit-for-purpose embedding selection rather than a universal model ranking.

This paper addresses that practical model-selection problem by presenting a reusable local benchmarking pipeline for embedding-model evaluation under corpus-specific RAG requirements. The pipeline transforms local document collections into controlled benchmark artifacts, evaluates candidate embedding models under one fixed protocol, and generates structured reports for retrieval effectiveness, operational behavior, cross-dataset consistency, and statistical decision support. Its intended output is decision support for practitioners: a way to compare candidate models transparently, preserve the underlying raw evidence, and regenerate dataset-level or aggregate reports when the evaluation context changes. \citep{GitHub_RagRetrievalBench}

The pipeline is demonstrated through a fixed experimental snapshot covering three document collections and five embedding models. The snapshot contains 60 source documents, 1{,}914 chunks, 150 synthetic queries, and 952 positive query--chunk relevance judgments (qrels, the ground-truth labels that mark which chunks are relevant to each query), providing a compact but heterogeneous test bed for examining language sensitivity, domain sensitivity, and operational tradeoffs. The broader questions of language and domain sensitivity are instantiated through English and Turkish education corpora and through Turkish education and law/regulation corpora, while operational decision support is instantiated through shared retrieval-quality, timing, consistency, and statistical outputs. Setup details, including provenance and controlled parameters, are provided in Section~\ref{experiments}; quantitative comparative evidence is reported in Section~\ref{results}.

\paragraph{Contributions and Research Questions}\label{contributions}
The paper makes four contributions to local RAG benchmarking practice. First, it defines a reusable pipeline for corpus-specific embedding-model selection under a controlled evaluation frame. Second, it integrates automatic report generation that consolidates retrieval quality, operational measurements, consistency diagnostics, and statistical outputs into decision-support artifacts. Third, it supports cross-dataset comparability through shared protocol execution and aggregate reporting. Fourth, it provides a publicly accessible implementation and benchmark artifacts to improve reuse and reproducibility in practitioner workflows. \citep{GitHub_RagRetrievalBench}

\paragraph{Research Questions}\label{research-questions}
Guided by these contributions, the empirical analysis addresses four broad research questions:

\begin{itemize}
\tightlist
\item
  RQ1 (Pipeline utility): Can the pipeline produce decision-support evidence---a ranked model shortlist with retrieval-quality metrics, operational-cost measurements, and query-level uncertainty estimates---for embedding-model selection on a local document collection?
\item
  RQ2 (Language sensitivity): Under one fixed benchmark protocol, can the pipeline expose language-conditioned ranking changes within this bounded snapshot?
\item
  RQ3 (Domain sensitivity): Under one fixed benchmark protocol, can the pipeline expose domain-conditioned ranking changes within this bounded snapshot?
\item
  RQ4 (Operational decision support): How does the pipeline combine retrieval effectiveness, operational cost, and query-level uncertainty into fit-for-purpose deployment guidance?
\end{itemize}

To make these questions answerable rather than rhetorical, we fix the operational meaning of the terms used throughout the paper. \emph{Decision-support evidence} (used interchangeably with \emph{actionable evidence}) is a per-model summary---retrieval-quality metrics, query latency and indexing throughput, and query-level uncertainty intervals---that is sufficient to choose a model under a stated deployment priority such as quality-first or latency-first. \emph{Pipeline utility} (RQ1) is the pipeline's ability to produce such a summary reproducibly from a local corpus while preserving the underlying raw artifacts. \emph{Consistency diagnostics} are the cross-condition stability summaries defined in Section~\ref{results} (all-metric coefficient of variation, worst-case metric drop, and score range). \emph{Statistical decision support} is the query-level bootstrap confidence-interval layer that indicates whether an observed difference is distinguishable from sampling noise, and it replaces inferential tests computed over only three dataset-level means.

These research questions are intentionally framed at the level of local benchmarking practice, but the empirical claims in this paper are bounded to the fixed experimental snapshot reported here. The evaluated evidence comes from three curated document collections, two languages, two domain settings, five embedding models, and one controlled retrieval configuration. Therefore, the paper does not claim universal model superiority. Instead, it evaluates whether the proposed workflow can produce reliable, reusable, and decision-relevant evidence for fit-for-purpose embedding-model selection under controlled local conditions.

Section~\ref{related-work} positions the paper against prior work, Section~\ref{method} describes the proposed workflow, Section~\ref{experiments} documents the experimental setup, Section~\ref{results} reports comparative evidence, and Section~\ref{conclusion} concludes with implications and future directions.
